\section*{Capítulo \ref{ch:b2:deep-time-climate} --- O clima ao longo do tempo geológico}

\begin{solution}{exo:b2:deep-time-climate:1}
(a) $900 + 37100 + 700 + 3 = \qty{38703}{PgC}$, cerca de 65 vezes os \qty{591}{PgC} do ar.
(b) $\num{3.5e6}\times0.012011 = \qty{42000}{PgC}$. (c) $591/42000 = \qty{1.4}{\percent}$.
\end{solution}

\begin{solution}{exo:b2:deep-time-climate:2}
(a) Entradas $11.3 + 3 + 12.3 + 8.1 + 2.0 = \qty{36.7}{Tmol.yr^{-1}}$; saídas $26 + 16 =
\qty{42}{Tmol.yr^{-1}}$, iguais dentro de suas incertezas. (b) $(11.3 + 3)/36.7 = 0.39$: bem
menos da metade; a maior parte da entrada é carbono sedimentar antigo. (c) $\num{3.5e6}/42 = \qty{8.3e4}{yr}$.
\end{solution}

\begin{solution}{exo:b2:deep-time-climate:3}
(a) \ce{Mg2SiO4 + 4 CO2 + 2 H2O -> 2 Mg^{2+} + 4 HCO3^- + SiO2}. (b) Com \ce{2 Mg^{2+} + 4 HCO3^-
-> 2 MgCO3 + 2 CO2 + 2 H2O}, a reação líquida é \ce{Mg2SiO4 + 2 CO2 -> 2 MgCO3 + SiO2}: dois mols
de dióxido de carbono por mol de forsterita. (c) $\qty{1}{PgC} = 10^{15}/12.011 =
\qty{8.33e13}{mol}$ de carbono, que exigem \qty{4.16e13}{mol} de forsterita, $4.16\times10^{13}
\times\qty{140.7}{g} = \qty{5.86}{Pg}$, quase seis bilhões de toneladas.
\end{solution}

\begin{solution}{exo:b2:deep-time-climate:4}
Glaciação hirnantiana, 445 a \qty{443}{Ma}, refrigerador; maior extensão do gelo neopaleozoico, 307 a
\qty{299}{Ma}, refrigerador; Turoniano, 94 a \qty{89}{Ma}, estufa (o mais quente do Fanerozoico);
PETM, 55.9 a \qty{55.7}{Ma}, estufa; EECO, 53 a \qty{49}{Ma}, estufa; início do manto de gelo
antártico, cerca de \qty{34}{Ma}, começo do refrigerador atual; mantos de gelo do norte,
\qty{2.7}{Ma}, refrigerador.
\end{solution}

\begin{solution}{exo:b2:deep-time-climate:5}
(a) $63000/(8.314\times288^2) = \qty{0.0914}{K^{-1}}$, que arredonda para 0.09. (b) $74000/(8.314
\times288^2) = \qty{0.107}{K^{-1}}$. (c) $\ln1.2/0.09 = \qty{2.03}{K}$ e $\ln 1.2/0.107 =
\qty{1.70}{K}$. (d) No estado estacionário o intemperismo deve igualar a desgaseificação, o que fixa $e^{k
\Delta T}$ e portanto $\Delta T$; a sensibilidade só determina quanto dióxido de carbono produz esse
aquecimento.
\end{solution}

\begin{solution}{exo:b2:deep-time-climate:6}
Aqui $k_T = 0.0914$. (a) Exato $\exp[7577.5\times(1/288 - 1/293)] = 1.567$; aproximação
$e^{0.0914\times5} = 1.579$: \qty{0.8}{\percent} acima. (b) Exato $e^{2.482} = 11.97$;
aproximação $e^{2.741} = 15.50$: \qty{30}{\percent} acima. (c) $\ln(\text{razão}) = (E_{\text
a}/R)[x/T_0^2 - x/(T_0(T_0 + x))] = (E_{\text a}/R)\,x^2/(T_0^2(T_0 + x))$. Igualando-o a $\ln1.1$:
$x^2 = 1.043\,(288 + x)$, logo $x = \qty{17.9}{K}$. Além de cerca de \qty{18}{K} de aquecimento a
forma exponencial superestima o intemperismo em mais de um décimo.
\end{solution}

\begin{solution}{exo:b2:deep-time-climate:7}
(a) $1444/276 = 5.23$, e $3.93\log_2 5.23 = \qty{9.4}{W.m^{-2}}$. (b) Quase exatamente os
\qty{9}{W.m^{-2}} do Sol: os dois forçamentos se compensam. (c) $(1444 - 276)/419.5 =
\qty{2.8}{ppm}$ por milhão de anos. Os \qty{3.4}{ppm} vêm de uma reta ajustada a todos os
dados, que pondera cada estimativa em vez das duas extremidades da curva suavizada.
\end{solution}

\begin{solution}{exo:b2:deep-time-climate:8}
(a) De $3000/20000 = \qty{0.15}{PgC.yr^{-1}}$ a $7000/3000 = \qty{2.3}{PgC.yr^{-1}}$, uma faixa
que contém os 0.3 a 1.5 do IPCC. (b) $35\times12.011/44.01 = \qty{9.6}{PgC.yr^{-1}}$, 6 a
32 vezes os 0.3 a \qty{1.5}{PgC.yr^{-1}} do PETM. (c) $3.93\log_2(2000/900) = \qty{4.5}{W.m^{-2}}$,
e $3\log_2 2.22 = \qty{3.5}{\celsius}$: cerca de setenta por cento dos \qty{5}{\celsius}. Outras
retroalimentações ou outra fonte de aquecimento, como o metano, devem completar o resto.
\end{solution}

\begin{solution}{exo:b2:deep-time-climate:9}
(a) Agora $1.2\,W_0e^{k\Delta T} = V = W_0$, logo $\Delta T = -\ln 1.2/0.09 = \qty{-2.03}{K}$.
(b) $278\times2^{-2.03/3} = \qty{174}{ppm}$. (c) Não: em qualquer estado estacionário o intemperismo iguala a
desgaseificação, que não mudou. A intemperabilidade maior é compensada por um clima mais frio.
\end{solution}

\begin{solution}{exo:b2:deep-time-climate:10}
(a) $k = 0.119$: $\tau = 591\times0.693/(0.08\times0.119\times3\times0.138) = \qty{104}{kyr}$.
(b) $S = 8$: $\tau = \qty{51.5}{kyr}$. (c) Ambas tornam o termostato mais rápido, pois $\tau \propto
1/(kS)$. $E_{\text a}$ é incerta por um fator de cerca de dois, $S$ nessas escalas de tempo por um
fator semelhante (3 contra 5 a \qty{8}{K}): ambas importam mais ou menos igualmente.
\end{solution}

\begin{solution}{exo:b2:deep-time-climate:11}
(a) $0.08\times5000 = \qty{400}{PgC}$ permanecem no ar: $C/C_0 = 991/591 = 1.677$, ou
\qty{466}{ppm}; $\Delta T = 3\log_2 1.677 = \qty{2.24}{K}$. (b) $W - W_0 = 0.138\,(e^{0.09\times2.24}
- 1) = \qty{0.0308}{PgC.yr^{-1}}$, \qty{22}{\percent} mais intemperismo; $5000/0.0308 =
\qty{162}{kyr}$. (c) $\tau = \qty{137}{kyr}$ e $\tau\ln10 = \qty{315}{kyr}$. (d) O aquecimento cresce
apenas como o logaritmo do dióxido de carbono, de modo que um grande excesso eleva o intemperismo menos do que a
lei linear supõe; a remoção começa mais devagar.
\end{solution}

\begin{solution}{exo:b2:deep-time-climate:12}
(a) Perto de $M^*$, $\mathrm d\Delta T/\mathrm dM = Sf/(A_0\ln 2\,(1 + fM^*/A_0)) = Sf/(A_0\ln 2\,
C^*/C_0)$ e $\mathrm dW/\mathrm dM = kW\,\mathrm d\Delta T/\mathrm dM$, com $W = V$ ali; o
inverso é $\tau'$. (b) $C^*/C_0 = 2^{2.03/3} = 1.597$: $\tau' = 137.4\times1.597/1.2 =
\qty{183}{kyr}$. (c) Mais lento: com muito dióxido de carbono cada petagrama de carbono acrescentado aquece menos,
porque o forçamento é logarítmico, de modo que o intemperismo responde menos a ele.
\end{solution}

\begin{solution}{pb:b2:deep-time-climate:1}
\textbf{1.} $1.3 + 1.3 + 4.5 + 4.2 = \qty{11.3}{Tmol.yr^{-1}}$, $\qty{0.136}{PgC.yr^{-1}}$.
\textbf{2.} $11.5\times0.012011 = \qty{0.138}{PgC.yr^{-1}}$.
\textbf{3.} Concordam a cerca de \qty{2}{\percent}: o intemperismo dos silicatos equilibra a desgaseificação
vulcânica, como exige o \omterm{def:b2:deep-time-climate:carbon-cycle}{ciclo do carbono de longo prazo}.
\textbf{4.} $\num{5.1e18}/0.02897 = \qty{1.76e20}{mol}$ de ar; $\num{591e15}/12.011 =
\qty{4.92e16}{mol}$ de carbono; $\num{4.92e16}/\num{1.76e20} = \num{2.80e-4}$, \qty{280}{ppm}.
\textbf{5.} $591/0.136 = \qty{4350}{yr}$.
\textbf{6.} $1 + 0.08 = 1.08$: apenas \qty{8}{\percent} a mais.
\textbf{7.} $63000/(8.314\times288^2) = \qty{0.0914}{K^{-1}}$, que arredonda para 0.09.
\textbf{8.} $e^{0.09} = 1.09$: \qty{9}{\percent} a mais por kelvin.
\textbf{9.} $\Delta T = S\log_2(1 + fM/A_0)$; sua derivada em 0 é $Sf/(A_0\ln2) =
3\times0.08/(591\times0.693) = \qty{5.86e-4}{K.PgC^{-1}}$.
\textbf{10.} $W_0k\times\num{5.86e-4} = 0.138\times0.09\times\num{5.86e-4} = \qty{7.28e-6}{yr^{-1}}$.
\textbf{11.} Seu inverso é $A_0\ln2/(fkSW_0)$, o $\tau$ da proposição, pois a
equação linearizada é $\mathrm dM/\mathrm dt = -(\mathrm dW/\mathrm dM)\,M$.
\textbf{12.} Os quatro: $\tau$ é inversamente proporcional a cada um.
\textbf{13.} $1/\num{7.28e-6} = \qty{1.37e5}{yr}$.
\textbf{14.} $137\times0.09/0.119 = \qty{104}{kyr}$.
\textbf{15.} $k_T = 0.107$: $137\times0.09/0.107 = \qty{115}{kyr}$.
\textbf{16.} Mais curtos por um fator de 1.5 a 7: o modelo de uma caixa e seu $f$ constante tornam o
termostato rápido demais.
\textbf{17.} $\num{50e6}/\num{1.37e5} = 364$: o termostato teve tempo de equilibrar o ciclo em cada
etapa, de modo que o declínio seguiu mudanças lentas da desgaseificação ou da intemperabilidade.
\textbf{18.} $16\times0.137 = \qty{2.2}{ppm}$.
\textbf{19.} $0.08\times5000 = \qty{400}{PgC}$; $278\times991/591 = \qty{466}{ppm}$.
\textbf{20.} $3\log_2 1.677 = \qty{2.24}{K}$.
\textbf{21.} $0.138\,(e^{0.202} - 1) = \qty{0.0308}{PgC.yr^{-1}}$, \qty{22.3}{\percent}.
\textbf{22.} $137\times\ln10 = \qty{315}{kyr}$.
\textbf{23.} O intemperismo pode ter sido mais sensível (escoamento, uma energia de ativação maior, as
retroalimentações lentas de uma sensibilidade do sistema terrestre maior); e o carbono orgânico em terra e no mar pode ter
crescido de novo e retirado carbono do ar mais depressa que o intemperismo dos silicatos.
\textbf{24.} Neste modelo o \omterm{def:b2:deep-time-climate:thermostat}{termostato do intemperismo dos silicatos} remove um excesso de carbono com um
tempo de decaimento de cerca de \textbf{\qty[text-series-to-math=true]{137}{} mil anos}.
\end{solution}
